% Vietnamese translation for OI Public Library
% Source-File: IOI中国国家候选队论文/2004/贝小辉.ppt
\documentclass[11pt]{article}
\usepackage{oipl-vn}

\newcommand{\Oh}{\mathcal{O}}

\title{Phân tích bài toán phân hoạch cây\\{\large Ghi chú theo slide}}
\author{Bei Xiaohui}
\date{IOI 2004}

\begin{document}
\maketitle

\origtitle{Tree partitioning problems}
\sourcepath{IOI 2004 / Bei Xiaohui.ppt}

\section{Bài toán gốc}

Các slide nhắc bài \emph{Strawberry} của NOI 2003, với các bộ kiểm thử
\texttt{test6} và \texttt{test9}. Trên cây, cần chọn hoặc chia các phần sao
cho giá trị trọng số đạt điều kiện tối ưu. Ký hiệu \(\operatorname{sum}(i)\)
và
\[
  x=\min\{\operatorname{sum}(i)\mid 1\le i\le k\}
\]
cho thấy tiêu chí phụ thuộc vào các tổng trên những phần được chia.

\section{Thuật toán dịch chuyển}

Các mốc \(W_{\min}\), \(W_{\mathrm{now}}\), \(W_{\min}(A)\) và
\(W_{\min}(Q)\) trên slide gắn với thuật toán dịch chuyển: bắt đầu từ một
phân hoạch, ta điều chỉnh biên giữa các phần của cây để cải thiện trọng số
xấu nhất. Khi chứng minh, cần xác định quan hệ ``phía trên'' giữa các phần
và chỉ ra mỗi lần dịch chuyển đều tiến về cấu hình tốt hơn hoặc giữ được
điều kiện cần.

\section{Độ phức tạp}

Slide nhắc dạng \(\Oh(k^2\operatorname{rd}(T)+kn)\), trong đó
\(\operatorname{rd}(T)\) là tham số phụ thuộc cấu trúc cây. Bài viết PDF còn
trình bày cách dùng thuật toán chuyển hóa ban đầu để tối ưu thuật toán dịch
chuyển, giảm số trường hợp cần xét.

\section{Ghi nhớ}

Với bài phân hoạch cây, điều khó thường nằm ở biên giữa các phần. Mô hình tốt
phải mô tả được cách biên này di chuyển và ảnh hưởng tới giá trị cực trị của
từng phần.

\end{document}
